\chapter*{Remerciements}
\addcontentsline{toc}{chapter}{Remerciements}

Je rends grâce à Dieu pour la force, la persévérance et l'intelligence accordées durant la réalisation de ce travail.

J'exprime ma profonde reconnaissance à l'Université de Kinshasa, à la Faculté Polytechnique et au Département de Génie Électrique et Informatique pour la formation reçue et le cadre scientifique offert.

Je remercie mon directeur de mémoire, le Prof.Dr.Ir. KYANDOGHERE KYAMAKYA Jean Vianney, ainsi que mes co-directeurs, le Dr. KAMBALE WITESYAVWIRWA Vianney et l'Ir.civ. MUGISHA MUHIRE Jean-Pierre, pour l'encadrement scientifique, les orientations méthodologiques et les remarques ayant permis d'améliorer ce travail.

J'adresse une reconnaissance particulière à mes frères et sœurs, BITATI MUKANDEZI Naomie Olive, MARHEGANE TSHIKONGA Ruth, MARHEGANE MPOYI Michel et BITATI KADIMASHI Benjamin, pour leur présence, leur soutien moral et leurs encouragements constants.

Je remercie affectueusement ma nièce, MAKENGO BITATI Meghan-Arielle, dont la présence a été pour moi une source de joie et de motivation.

Mes remerciements vont également à mes amis SITU SANZU Aristote et KAPOLA Joseph, pour leur disponibilité, leurs encouragements et les moments de partage qui ont accompagné ce parcours.

Je tiens aussi à exprimer ma gratitude à mon oncle, le Dr.Ir.MUTEBA TSHIASUMA Adolphe, pour ses conseils, son soutien et son intérêt constant pour l'aboutissement de mon parcours académique.

Je remercie enfin, de manière générale, tous mes collègues de promotion et tous ceux avec qui j'ai partagé les échanges académiques, les travaux pratiques, les discussions techniques et les efforts quotidiens qui ont marqué cette formation. Que toutes les personnes qui ont contribué, directement ou indirectement, à la réalisation de ce mémoire trouvent ici l'expression de ma sincère gratitude.
